\section{A subset-sum precode above the binary Gilbert--Varshamov benchmark}
\label{app:subset-sum-gv}

The support catalog of the precode is another design parameter. The
Newton precode $B_8$ is MDS, so it maximizes local distance, but it
permits many distinct supports. The following smaller circuit trades
local distance for fewer supports and gives a stronger complete distance
bound with a smaller failure margin.

Index eight inputs by $j\in\F_2^3$ and define
\[
 y_i=\sum_{j:\langle i,j\rangle=0}x_j,\qquad
 B_{\rm sub}(x)=(x_0,\ldots,x_7,y_0,\ldots,y_7).
\]
The binary dot product selects the subsets; the sums are evaluated over
$\F_p$. Thus $y_0$ sums all eight inputs and each other $y_i$ sums four.
Every entry of the matrix of $B_{\rm sub}$ is zero or one. Let $E$ be
the instance of~\eqref{eq:width-family} with $w=1$, $q=1$, and $L=2$,
and with $B_8$ replaced by the subset-sum precode $B_{\rm sub}$ applied
blockwise. Thus $E$ applies $B_{\rm sub}$, then two independently
permuted plain accumulators.

\begin{theorem}\label{thm:subset-sum-gv}
For $p=2^{128}-159$, $k=2^{20}$, and $n=2k$, the encoder $E$ has rate
$1/2$ for every setup. Except with probability less than $2^{-22.05}$
over the two independent uniform permutations, its minimum distance is
at least $263403$, which exceeds $12.56\%$ of $n$. Encoding uses exactly
$3.3125n-2=6946814$ field additions, with no subtractions or field
multiplications.
\end{theorem}

This relative distance exceeds the binary rate-$1/2$
Gilbert--Varshamov (GV) benchmark $H_2^{-1}(1/2)\approx0.1100278644$.
The theorem concerns the specified finite prime-field code. It does not
compare its actual distance with that of the same circuit over $\F_2$,
or give an asymptotic family theorem.

\paragraph{Circuit.}
The systematic copy gives injectivity. To compute the remaining outputs,
first form $x_0+x_j$ for $j\in\{1,2,3\}$, using the usual three-bit
representation of the indices. Together with $0$, these three indices
form a two-dimensional subspace of $\F_2^3$, so every nonzero $i$ is
orthogonal to at least one of them. For such a $j$, the shared sum
$x_0+x_j$ and two further inputs compute $y_i$ in two additions.
Finally, $y_0=y_1+((x_1+x_3)+(x_5+x_7))$ costs four additions. The
block therefore uses $3+14+4=21$ additions, and the two scans use
$2(n-1)$.

\paragraph{Support catalog.}
Let $Z$ be a set of local output coordinates constrained to zero, and let
$V_Z$ be the kernel of those rows of $B_{\rm sub}$. An exact zero set
must be closed under row span: every local row in the span of $Z$ also
vanishes. Conversely, if $Z$ is closed and has rank $j<8$, its kernel
has dimension $8-j$ and each remaining coordinate is a nonzero linear
functional on that kernel. Since $p>16$, their proper hyperplanes cannot
cover the kernel. Hence every proper closed set is realized as an exact
zero set.

Enumerating the $2^{16}$ subsets by exact elimination over $\F_p$ gives
the polynomial $S(z,\tau)$ in Table~\ref{tab:subset-source}. Its coefficient
of $z^s\tau^d$ counts supports of size $s$ whose shortened source space
has dimension $d$. The constant term represents the zero block.

\begin{table}[t]
\centering
\caption{Exact shortening catalog of the eight-input subset-sum block.}
\label{tab:subset-source}
\begin{tabular}{rl@{\qquad}rl}
\toprule
$s$ & $[z^s]S(z,\tau)$ & $s$ & $[z^s]S(z,\tau)$\\
\midrule
0 & $1$ & 10 & $3045\tau^2+154\tau^3$\\
5 & $7\tau$ & 11 & $3514\tau^3+7\tau^4$\\
6 & $77\tau$ & 12 & $1785\tau^4$\\
7 & $28\tau$ & 13 & $560\tau^5$\\
8 & $154\tau+84\tau^2$ & 14 & $120\tau^6$\\
9 & $772\tau+511\tau^2$ & 15 & $16\tau^7$\\
 & & 16 & $\tau^8$\\
\bottomrule
\end{tabular}
\end{table}

\paragraph{Finite distance proof.}
For $n/16$ independent blocks, the dimension-weighted source moment at
support $s$ is $[z^s]S(z,\tau)^{n/16}$. The projective-pair and
joint-renewal bounds of Appendices~\ref{app:fresh-pairs}
and~\ref{app:joint-renewal} are uniform in the source forms and their
successive quotients. They therefore apply with this source polynomial;
their scan inequalities and finite sampling corrections are unchanged.

These two bounds alone prove that $E$ has distance at least $239076$
($>11.4\%$ of $n$), except with probability less than $2^{-24.26}$.
Take forbidden weight $e=239075$. The sparse bound of
Appendix~\ref{app:all-one}, with source counts from $S(z,1)^{n/16}$,
covers $1\le s\le31$. The remaining domain is
$32\le s\le4e$ and $\lceil s/2\rceil\le h\le2e$. It is partitioned
recursively into rectangles, and each cell is a leaf: 62 leaves use
the projective-pair bound and 61 use the joint-renewal bound. The
verifier recomputes the local catalog, checks every partition and
vertex bound, and checks each joint-renewal inequality over the entire
pool simplex. Adding all sparse and dense terms with upward rounding
gives a failure probability with base-two logarithm at most
$-24.265906406170$.

For the stronger forbidden weight $e=263402$ of
Theorem~\ref{thm:subset-sum-gv}, use the complete-run bound
of Appendix~\ref{app:complete-runs} with the same catalog $S$.
The sparse support count again covers $1\le s\le31$. The dense cover
has 65 projective-pair leaves and 131 complete-run leaves, covering
every remaining feasible profile. The total failure probability has
base-two logarithm at most $-22.051928329751$, which proves the theorem.
The verifier reconstructs the exact local catalog. Both certificates
and their verifiers are listed in Appendix~\ref{app:reproduction}.

The subset-sum encoder thus has a larger proved distance and fewer
additions than the Newton encoder of Appendix~\ref{app:complete-runs},
whose distance exceeds $11.7\%$ of $n$ with failure probability below
$2^{-49.64}$, but a smaller failure margin. We did not time this variant.
